\section{Experiments}
\label{sec:experiments}


\subsection{Experimental Setup}
\label{sec:setup}

\textbf{Datasets.}
\label{sec:datasets}
We evaluate our framework on the \textit{DiDi GAIA Open Dataset}. The dataset contains high-resolution GPS trajectories and trip dispatch logs from November 2016, covering approximately 7.06 million completed trips and over 1.09 billion GPS pings. We use 9 November 2016, a representative weekday, as the primary evaluation instance, which includes 36,799 workers and 224,219 tasks in the core urban area around Chengdu. Each unique driver is modeled as a worker, with their first GPS ping used to determine the online time and initial location, and an 8-hour shift duration is applied. Each completed trip is converted into one task, where the dispatch timestamp is used as the release time and the billing-completion timestamp as the expiry deadline. Following the methodology in \cite{CR-04}, we compute task revenue $m_j$ using a base fare $\beta_{base} = \$2.00$ and a distance rate $\beta_{dist} = \$1.50$/km, where $\alpha_j$ is the Manhattan distance between $l_j^o$ and $l_j^d$. Total platform revenue is the aggregate $m_j$ of all completed tasks. 

To test whether the observed strategy behavior generalizes beyond dense ride-hailing scenarios, we also evaluate utilising the \textit{Gowalla} location-based social network dataset. We focus on the Austin, Texas region during September 2010, which provides the densest usable check-in cluster in the dataset. Since raw LBSN check-ins are much sparser than ride-hailing requests, we apply temporal compression by projecting all check-ins onto a single 24-hour reference window while preserving their wall-clock times and spatial distribution. Each unique user-day pair is treated as one worker, and each check-in is converted into one task with a 30-minute expiry window. Because Gowalla does not contain trip destinations or fares, drop-off locations and revenues are generated using the distance-proportional proxies described above. After compression and worker downsampling, the Gowalla evaluation instance contains 8,758 workers and 43,788 tasks.

\textbf{Implementation Details.}
\label{sec:framework}
We implement a dataset-agnostic discrete event simulator (DES) driven by a min-heap priority queue ordered by timestamps, ensuring that all events are processed chronologically without artificial batching or future lookahead, consistent with the instantaneous matching constraint in online spatial crowdsourcing~\cite{tong2019survey}. The simulator supports worker release, task release, task completion, and task expiration events. Each assignment strategy is implemented through a pair of handlers, $(\textsc{NewTask}, \textsc{FreeWorker})$, which receive the current system state and return at most one assignment per event. Individual worker earnings are tracked as the sum of revenues $m_j$ from their completed tasks. Available workers are maintained in an incrementally updated $k$-d tree spatial index based on their current locations, enabling efficient $k$-nearest-neighbor queries without rebuilding the full index at every event. Distances are computed using Manhattan distance under a locally-scaled flat-earth projection, where longitude is rescaled by $\cos(\bar{\phi})$ at the city's mean latitude $\bar{\phi}$. This approach reduces geodetic error to below 0.15\% over the urban extent while preserving the $\mathcal{O}(k \log n)$ query efficiency required for city-scale dispatch. Travel times assume a constant average urban speed of 30\,km/h. Each candidate pair is filtered by two feasibility constraints: the worker must reach the pickup location before the task deadline, and must complete the trip before their shift deadline. Candidate assignments violating either condition are rejected before scoring, which prevents infeasible matches without introducing additional penalty terms.


\textbf{Baselines.}
We compare our approach with representative online, fairness-aware, randomized, and batch-based assignment baselines. Greedy minimizes pickup distance only by setting $\lambda_1,\lambda_2=0$, while Least Allocated Worker First (LAF) assigns tasks to workers with the fewest completed tasks as a pure fairness baseline. EWMA-Only uses only the proposed EWMA idle-time fairness signal without utility or starvation terms, and Fixed-Composite applies a static grid-searched weight configuration of our composite framework. BiRanking (BRK) is a randomized ranking-based policy inspired by \cite{[CITATION]}, which greedily matches feasible workers and tasks according to fixed random ranks. FATP-ANN is an existing fairness-aware online assignment method \cite{CR-08}. TSGF Sampling adapts the Two-Sided Group Fairness framework \cite{[CITATION]} by randomly selecting operator-profit, worker-fairness, or task-fairness heuristics at each dispatch event. Cost-Balancing (CB) \cite{[CITATION]} and MMD-Batch (MMD-BB) \cite{CR-11} are online batching baselines that defer decisions to balance matching cost or minimize worst-case delay. Discrete Review LP \cite{[CITATION]} periodically solves a global bipartite assignment problem via the Hungarian algorithm, while ONRTA-RT and ONRTA-OP \cite{[CITATION]} represent randomized-threshold and optimal-threshold variants of the Online Non-Rejection Task Assignment policy.


\textbf{Evaluation Metrics.}
To evaluate how well the proposed composite assignment framework navigates the tripartite trade-off among efficiency, fairness, and starvation prevention, we report metrics covering platform performance, worker equity, service reliability, and computational cost. Platform efficiency is measured by the \textit{Task Assignment Ratio} (TAR), defined as the fraction of successfully matched requests over all task arrivals ($|M|/|\mathcal{T}|$), together with \textit{Total Platform Revenue} ($k$), which reflects the economic viability of the assignment strategy. Supply-side equity is evaluated using Jain's Fairness Index (JFI) \cite{JFI-Original}, computed over the vector of cumulative task counts $\mathbf{x}$ across active workers:
$JFI(\mathbf{x}) = \left.{(\sum_{i=1}^{n} x_i)^2} \middle/n {\sum_{i=1}^{n} x_i^2} \right.$,
where $n$ is the number of workers and a value closer to 1.0 indicates a more equitable task distribution. Demand-side reliability is captured by \textit{Mean Wait Time}, which measures the average interval from task release $r_j$ to worker pickup, and \textit{Average Pickup Distance}, which measures the mean empty travel distance required for workers to reach pickup locations. Together, these metrics indicate whether a strategy achieves responsiveness through efficient spatial matching or at the cost of longer detours. Finally, we report \textit{Peak Backlog}, defined as the maximum number of simultaneously unassigned tasks during a simulation day, to assess resilience to demand spikes, and \textit{Wall-clock Runtime} to evaluate whether the method remains computationally practical for city-scale deployment.

\subsection{Effectiveness}
\label{sec:overall_results}
\begin{table}[t]
\scriptsize
\renewcommand\arraystretch{1}
\setlength\tabcolsep{2pt}
\caption{Strategy comparison on the DiDi Chengdu held-out test day
  (9 Nov 2016; 36,799 workers, 224,219 tasks).
  $^\dagger$ proposed method. Bold = column-best among strategies with TAR $\geq$~0.85.
  $^\ddagger$ FATP-ANN exhibits degenerate throughput at city scale (TAR $=$ 0.599,
  peak backlog $=$ 39{,}298), attributable to threshold-parameter sensitivity to
  distribution shift; values shown in parentheses and excluded from comparison.}
\label{tab:didi_results}
\begin{tabular}{@{}l cccc c c@{}}
\toprule
\textbf{Strategy}
  & \textbf{TAR} $\uparrow$
  & \textbf{Wait (m)} $\downarrow$
  & \textbf{JFI (tasks)} $\uparrow$
  & \textbf{JFI (earn.)} $\uparrow$
  & \textbf{Rev.\ (k\$)} $\uparrow$
  & \textbf{Time (s)} $\downarrow$ \\
\midrule
Greedy (baseline)
  & 0.863 & 2.62 & 0.594 & 0.613 & 2{,}765.4 & 357 \\
k-NLF ($k{=}15$)$^\dagger$
  & \textbf{0.863} & 3.45 & 0.649 & 0.669 & \textbf{2{,}765.8} & \textbf{38} \\
Composite (static)$^\dagger$
  & \textbf{0.863} & 3.20 & 0.589 & 0.625 & \textbf{2{,}765.8} & 98 \\
\midrule
LTF
  & 0.863 & 9.63 & \textbf{0.705} & \textbf{0.700} & 2{,}765.1 & 319 \\
BiRanking (BRK)
  & 0.857 & 12.34 & 0.592 & 0.636 & 2{,}750.9 & 434 \\
FATP-ANN$^\ddagger$
  & (0.599) & (8.46) & (0.957) & (0.887) & (2{,}206.9) & 12{,}410 \\
ONRTA-RT
  & 0.855 & 12.17 & 0.643 & 0.657 & 2{,}749.8 & 416 \\
Disc.\ Review LP
  & 0.863 & \textbf{1.64} & 0.517 & 0.699 & 2{,}765.5 & 5{,}060 \\
\bottomrule
\end{tabular}
\end{table}

\begin{table}[t]
\scriptsize
\renewcommand\arraystretch{1}
\setlength\tabcolsep{2pt}
\caption{Strategy comparison on the Gowalla Austin dataset
  (Sep 2010, temporally compressed; 8,758 workers, 43,788 tasks, worker-to-task ratio~1:5).
  $^\dagger$ proposed method. Bold = column-best among strategies with TAR $\geq$~0.99.
  FATP-ANN achieves high JFI at a severe throughput cost (TAR $=$ 0.700 vs $\geq$~0.987 elsewhere);
  values shown in parentheses and excluded from bolding.
  Revenue is a synthetic distance-proportional proxy, as Gowalla contains no recorded fares.}
\label{tab:gowalla_results}
\begin{tabular}{@{}l cccc c c@{}}
\toprule
\textbf{Strategy}
  & \textbf{TAR} $\uparrow$
  & \textbf{Wait (m)} $\downarrow$
  & \textbf{JFI (tasks)} $\uparrow$
  & \textbf{JFI (earn.)} $\uparrow$
  & \textbf{Rev.\ (k\$)} $\uparrow$
  & \textbf{Time (s)} $\downarrow$ \\
\midrule
Greedy (baseline)
  & 0.998 & 5.12 & 0.591 & 0.621 & 297.5 & \textbf{4.4} \\
k-NLF ($k{=}15$)$^\dagger$
  & 0.998 & 5.16 & 0.672 & 0.693 & 297.5 & 27.3 \\
Composite (static)$^\dagger$
  & 0.998 & 5.11 & 0.612 & 0.654 & \textbf{297.6} & 28.1 \\
\midrule
LAF
  & 0.998 & 11.82 & \textbf{0.894} & \textbf{0.877} & \textbf{297.6} & 91.0 \\
BiRanking (BRK)
  & 0.988 & 15.87 & 0.360 & 0.579 & 294.5 & 115.7 \\
FATP-ANN
  & 0.700 & 9.33 & (0.950) & (0.921) & 227.1 & 209.8 \\
ONRTA-RT
  & 0.987 & 15.97 & 0.671 & 0.708 & 294.3 & 108.6 \\
Disc.\ Review LP
  & \textbf{0.999} & \textbf{2.72} & 0.575 & 0.601 & \textbf{297.6} & 142.9 \\
\bottomrule
\end{tabular}
\end{table}



\subsection{Efficiency \& Scalability}
\label{sec:scalability}

% TODO: insert scalability plots and prose once scalability\_tasks\_v2.csv arrives from cluster.

\subsection{Ablation \& Sensitivity Analysis}
\label{sec:ablation}

\subsubsection{Impact of Candidate Pool Size ($k$).}
Figure~\ref{fig:k_sweep} shows how varying the candidate pool size $k$ affects JFI, average wait time, and wall-clock runtime for both k-NLF and Composite.

\begin{figure}[t]
  \includegraphics[width=\textwidth]{figures/5.4/k_sweep.pdf}
  \caption{Impact of candidate pool size $k$ on (a) task-count fairness (JFI), (b) average task wait time, and (c) wall-clock runtime for k-NLF and Composite (DiDi, 9 Nov 2016). The dashed grey line indicates the Greedy baseline; the vertical guide marks the paper default $k{=}15$. Fairness returns diminish beyond $k{\approx}25$ while runtime remains stable, validating $k{=}15$ as a practical operating point.}
  \label{fig:k_sweep}
\end{figure}

The parameter $k$ determines the search radius for fairness interventions, governing the trade-off between local proximity and distributive equity. We sweep $k \in \{3, \dots, 100\}$ across both proposed strategies, using the Greedy strategy as a baseline anchor.

\paragraph{Fairness Returns and the Local Advantage.} 
As illustrated in Fig.~\ref{fig:k_sweep}(a), worker fairness (JFI) generally increases with $k$ for both strategies. The task-count equity of k-NLF exhibits high sensitivity to pool size, rising from 0.593 ($k=3$) to 0.731 ($k=100$). Crucially, k-NLF surpasses the fairness of the global LAF baseline reported in Table~\ref{tab:didi_results} (JFI = 0.705) once $k \geq 50$. This suggests that in high-density urban environments, localized selection from a sufficient pool can more effectively mitigate task-monopolization hotspots than global optimization. In contrast, the Composite strategy exhibits more conservative fairness growth, remaining substantially below k-NLF across all pool sizes, which reflects its focus on temporal participation equity rather than snapshot task counts.

\paragraph{The U-Shaped Efficiency Response.} 
Fig.~\ref{fig:k_sweep}(b) reveals a distinct \textit{U-shaped response} in average task wait times for both strategies. At low search radii ($k=3$), wait times are significantly higher ($\sim$4.6\,min) than the Greedy baseline. This occurs because a pool of three candidates is often too restrictive to find a geographically close worker who also satisfies the fairness criteria. As $k$ increases, wait times drop to a global minimum at $k=25$ (3.32\,min for k-NLF; 3.16\,min for Composite). This identifies a ``Pareto Sweet Spot'' around $k \approx 25$ where the pool is large enough to contain workers who are both relatively nearby and underserved. While $k=25$ minimizes latency, $k=15$ was selected as the paper default because the wait-time penalty relative to the minimum is marginal ($+0.11$\,min) while still providing robust fairness gains and a smaller computational footprint for high-frequency deployment.

\paragraph{Scalability and Overhead.} 
A primary claim of this work is that $O(k)$ dispatching enables city-scale fairness at a tractable cost. Fig.~\ref{fig:k_sweep}(c) demonstrates that after an initial drop in runtime as $k$ moves from 3 to 15---likely attributable to the increased ease of finding feasible candidates in a larger pool---the wall-clock runtimes remain remarkably stable for $k \geq 25$. While k-NLF and Composite incur higher overhead than the proximity-only Greedy baseline, their runtimes are significantly lower than the global LP and LAF models reported in Table~\ref{tab:didi_results}. This stability confirms that the search radius can be expanded to maximize fairness without triggering the superlinear scaling bottlenecks characteristic of globally-coupled dispatchers.

\subsubsection{Fairness Signal Comparison.}
\label{sec:signal_comparison}

\begin{table}[t]
\scriptsize
\renewcommand\arraystretch{1}
\setlength\tabcolsep{2pt}
\caption{Comparison of $k$-NN fairness signals at $k{=}15$ (DiDi, 9 Nov 2016).
  $^\dagger$ proposed method. Bold = column-best among fairness-signal strategies.
  CV = coefficient of variation of the named quantity across workers (lower $\downarrow$ = more uniform distribution).
  k-NTF-EPH directly optimises earnings equity (CV earn.); k-NTF-IR directly optimises idle-time equity (CV idle).
  k-NLF achieves the best value in every fairness column, including those each alternative directly optimises.}
\label{tab:signal_comparison}
\begin{tabular}{@{}l cccc c c@{}}
\toprule
\textbf{Signal / Strategy}
  & \textbf{JFI (tasks)} $\uparrow$
  & \textbf{JFI (earn.)} $\uparrow$
  & \textbf{CV (idle)} $\downarrow$
  & \textbf{CV (earn.)} $\downarrow$
  & \textbf{Wait (m)} $\downarrow$
  & \textbf{Time (s)} $\downarrow$ \\
\midrule
Greedy (baseline)
  & 0.562 & 0.619 & 1.418 & 0.949 & 3.40 & 31 \\
\midrule
k-NTF-EPH ($k{=}15$)
  & 0.613 & 0.645 & 1.560 & 0.863 & 3.25 & \textbf{35} \\
k-NTF-IR ($k{=}15$)
  & 0.618 & 0.651 & 1.563 & 0.860 & 3.34 & \textbf{35} \\
k-NLF ($k{=}15$)$^\dagger$
  & \textbf{0.647} & \textbf{0.669} & \textbf{1.467} & \textbf{0.824} & 3.48 & 133 \\
\midrule
Composite (static)$^\dagger$
  & 0.589 & 0.626 & 1.602 & 0.918 & \textbf{3.21} & 114 \\
\bottomrule
\end{tabular}
\end{table}

\subsubsection{Fairness Weight Sensitivity.}
\label{sec:fw_sensitivity}

Figure~\ref{fig:fw_sweep} reports the effect of the fairness weight $\lambda_f$ on worker participation equity and task tail latency for the Composite strategy, with all other hyperparameters held fixed at their paper defaults.

\begin{figure}[t]
  \includegraphics[width=\textwidth]{figures/5.4/fairness_weight_sweep.pdf}
  \caption{Sensitivity of the Composite dispatcher to the fairness weight $\lambda_f$ (all other parameters fixed: $\lambda_s{=}0$, $\lambda_u{=}1$, $\gamma{=}0.1$, $k{=}15$; DiDi, 9 Nov 2016). Panel~(a): JFI rate (worker participation equity). Panel~(b): P95 task wait time (tail latency). Dashed grey = Greedy baseline; dash-dot purple = k-NLF reference. The vertical guide marks the paper default $\lambda_f{=}1.6$.}
  \label{fig:fw_sweep}
\end{figure}

Figure~\ref{fig:fw_sweep} examines how varying the fairness weight
$\lambda_f \in [0, 3]$ affects worker participation equity and task tail
latency in the Composite dispatcher, with all other parameters fixed at
their paper defaults ($\lambda_s{=}0$, $\lambda_u{=}1$, $\gamma{=}0.1$,
$k{=}15$).

\paragraph{The Critical Role of the EWMA Signal.}
At $\lambda_f{=}0$, the Composite dispatcher reduces to a
proximity-only selector restricted to $k{=}15$ candidates. This
degenerate configuration produces a JFI rate of 0.852---\textit{below}
the Greedy baseline (0.864)---because the spatial restriction of
$k{=}15$ limits worker diversity without any compensating fairness
mechanism. This result directly motivates the EWMA signal: even a small
fairness weight of $\lambda_f{=}0.2$ is sufficient to cross the Greedy
baseline, lifting JFI rate to 0.870 in a single step.

\paragraph{Plateau Robustness.}
As illustrated in Fig.~\ref{fig:fw_sweep}(a), for $\lambda_f \geq 0.2$
the JFI rate stabilises in the range $[0.866, 0.871]$, remaining
consistently above the Greedy baseline across the entire sweep. The
peak of 0.871 occurs at $\lambda_f \approx 1.4$; the paper default of
$\lambda_f{=}1.6$ lies within this plateau at 0.867. This near-flat
response confirms that the Composite dispatcher is robust to precise
hyperparameter selection---operators do not need to fine-tune
$\lambda_f$ to achieve consistent fairness improvements.

\paragraph{Dual Benefit on Task Tail Latency.}
Fig.~\ref{fig:fw_sweep}(b) reveals that increasing $\lambda_f$ also
reduces P95 task wait time. At $\lambda_f{=}0$, P95 is 14.1\,min,
marginally exceeding the k-NLF reference (14.0\,min). As $\lambda_f$
increases, P95 drops to a sustained band of $12.9$--$13.1$\,min for
$\lambda_f \geq 0.6$, a reduction of approximately 1.0--1.2\,min at
the paper default. This improvement arises because the EWMA signal
increasingly redirects assignments away from repeatedly-served workers
toward underserved workers who are also spatially distributed, reducing
the tail of long-wait outliers. Notably, TAR remains flat across the
entire sweep ($\Delta\text{TAR} < 0.0001$), confirming that these
efficiency gains come at zero throughput cost.

\paragraph{Design Implication.}
The results collectively support a practical design guideline: any
$\lambda_f \in [0.2, 3.0]$ delivers consistent fairness and latency
improvements over both the pure-utility baseline ($\lambda_f{=}0$) and
the Greedy dispatcher. The paper default $\lambda_f{=}1.6$ was selected
from within this stable region; the near-identical performance across
the plateau means that grid-search precision is not required in
deployment, lowering the operational cost of adopting the Composite
framework.